\documentclass[answers]{exam}

\usepackage[margin=1in]{geometry}

\usepackage{comment}
\usepackage{graphicx}
\usepackage{hyperref}
\usepackage{listings}

\pagestyle{headandfoot}
\firstpageheader{CSCI 221}{}{Testing and Debugging Assignment}
\runningheader{CSCI 221}{}{Testing and Debugging Assignment}


\begin{document}
\title{Testing and Debugging Assignment (Evaluation)}
\author{CSCI 221: Computer Science Fundamentals 2}
\date{Spring 2026}
\maketitle


This assignment is designed to test your understanding of testing and debugging. 
Feel free to collaborate with others and use resources, but all code and writeups must be your own. 
%For each of the questions, you should also write code to test your results and a makefile that compiles your code. 
%You can use one makefile for all compilations, or you can have separate makefiles for different questions. 
For this assignment, you should submit your code along with a writeup that contains the answers to the written questions. 

\textbf{Due Date:}
Friday, February 20th at 1:00 pm. 

\begin{questions}

%\question[18]
%\textbf{Design.}
%header vs c
%names
%functions
%comments

\question[33]
\textbf{Testing.}
Using the method specified, write a series of test cases for the specified problem. 
For each test case, you should describe: 
\begin{enumerate}
\item The input to the test case. 
\item The expected behavior (output) of the test case. 
\item An explanation of how you chose the test case using the specified methodology. 
This does not need to be long, a sentence or two should be sufficient. 
\item What errors you expect the test case to catch. 
\end{enumerate}
\begin{parts}
\part[6]
For this part and the next, you will be considering the following problem:

This function will count how many times each day of the week appears in a range of calendar years.
The input will consist of one character representing the day of the week, and two unsigned integers representing the starting and ending year, respectively. 
(Both years should be included in the count.)
The characters used to represent the days of the week are \texttt{[S,M,T,W,h,F,a]} for Sunday through Saturday, respectively. 
The function only handles the years between 2000 and 2050. 

Create tests for this problem using equivalence testing. 

\begin{solution}
\begin{enumerate}
    \item Testing on Days of the week
    \begin{enumerate}
        \item Valid Cases: \begin{verbatim}
            
        \end{verbatim}
    \end{enumerate}
\end{enumerate}    
\end{solution}

\part[6]
Create tests for the problem from the previous part using boundary testing. 

\part[3]
For this part and the next two parts, you will be considering the following code:
\begin{verbatim}
unsigned int turn = 1;
unsigned int last_turn = 0;
unsigned int first_score = 0;
unsigned int second_score = 0;

void reset(void) {
    turn = 1;
    last_turn = 0;
    first_score = 0;
    second_score = 0;
}

unsigned int turn(int player, unsigned int dice) {
    if (player == turn) {
        if (player == 1) {
            for (int i=0; i < dice; i++) {
                first_score += die_roll();
            }
            if (first_score > 100) {
                turn = 0;
                return 2; 
            }
        }
        if (player == 2) {
            for (int i=0; i < dice; i++) {
                second_score += die_roll();
            }
            if (second_score > 100) {
                turn = 0;
                return 1;
            }
        }
        if (last_turn == 1) {
            turn = 0;
            if (first_score > second_score) {
                return 1;
            } else if (second_score > first_score) {
                return 2;
            } else {
                return 3;
            }
        }
        if (dice == 0) {
            last_turn = 1;
        }
        turn = 3 - turn;
    }
    return 0;
}
\end{verbatim}
This code plays the game Greed that students in CSCI 121 have to implement strategies for. 
If you want to see the rules to the game, look at the pdf file in the folder with the assignment. 

The code takes as input an integer representing which player is trying to take a turn, and an unsigned integer representing how many dice they want to roll. 
It tracks the state of the game, and returns an integer indicating the state at the end of the turn. 
If the game is not over, the function returns 0. 
If the game is over, then the function returns the winning player, represented by 1 or 2. 
In the event of a tie, the function returns 3. 

Do you need to write any additional code to test this code?
If so, what do you need to write?
If not, why not?

\part[9]
Create tests for the code from the previous part using control flow testing. 
You should submit an image of your control flow graph along with your test cases. 

\part[9]
Create tests for the code from the previous part using state-based testing. 
You should submit an image of your finite state machine along with your test cases. 
\end{parts}


\question[12]
\textbf{Debugging.}
Consider the following code, which attempts to compute exponentials for non-negative integers. 
\begin{verbatim}
unsigned int f(unsigned int x, unsigned int y) {
    int r = 1;
    int s = x;
    while (y > 1) {
        if (y % 2 == 1) {
            r = x * r;
        }
        s = s * s;
        y = y / 2;
    }
    return r * s;
}
\end{verbatim}
Your job is to debug the code so that it performs as specified. 
You should submit the following:
\begin{enumerate}
\item An explanation of any defects in the original code. 
\item For each defect, a brief description of how you found that defect. 
Your description should contain any experiments you did to verify your analysis. 
A few sentences should be sufficient to do this. 
\item The code with the defects fixed and working. 
\end{enumerate}
As part of the debugging process, modify the code so that it follows good style. 

\end{questions}

\end{document}